\documentclass[10pt,a4paper]{amsart}
\usepackage[margin=19mm]{geometry}
\usepackage{amsmath,amssymb,mathtools}
\usepackage[scr=rsfs,cal=euler]{mathalfa}
\usepackage{xfrac}
\usepackage[unicode,hidelinks,bookmarks=false]{hyperref}
\newcommand{\ie}{i.e.\ }
\newcommand{\cf}{cf.\ }
\newcommand{\resp}{respectively\ }
\newcommand{\Subsection}{\subsection}
\providecommand{\nobreakdash}{\nobreak\hbox{-}}
\setlength{\emergencystretch}{2em}
\multlinegap0pt
\usepackage{bm}
\usepackage{bbm}
\usepackage{dsfont}
\usepackage{xspace}
\usepackage{cancel}
\usepackage{mathtools}
\usepackage{cleveref}
\usepackage{amsthm}
\makeatletter
\@ifundefined{lemma}{\newtheorem{lemma}{Lemma}}{}
\@ifundefined{proposition}{\newtheorem{proposition}{Proposition}}{}
\makeatother
\DeclareMathOperator{\LD}{LD}
\DeclareMathOperator{\blocks}{blocks}
\DeclareMathOperator{\rem}{rem}
\newcommand{\Lcal}{\mathcal{L}}
\newcommand{\N}{\mathbb{N}}
\newcommand{\R}{\mathbb{R}}
\newcommand{\Z}{\mathbb{Z}}
\newcommand{\ba}{\mathbf{a}}
\newcommand{\bb}{\mathbf{b}}
\newcommand{\bu}{\mathbf{u}}
\newcommand{\conv}[1]{\operatorname{conv}\left(#1\right)}
\newcommand{\lp}{\left(}
\newcommand{\nvol}[1]{\operatorname{nvol}\left(#1\right)}
\newcommand{\rp}{\right)}
\newcommand{\spn}[1]{\operatorname{span}\left(#1\right)}
\newcommand{\st}{\, \mathrm{:}\,}
\title[On Lattice Diameter Segments and A Discrete Borsuk Partition]{On Lattice Diameter Segments and A Discrete Borsuk Partition Problem}
\author{Anouk E. Brose, Jes\'us A. De Loera, Gyivan Lopez-Campos, Antonio J. Torres}
\date{Source: arXiv:2508.20009v1 (CC BY 4.0)}
\begin{document}
\maketitle

\noindent\emph{Source and licence.} On Lattice Diameter Segments and A Discrete Borsuk Partition Problem. Version arXiv:2508.20009v1 (CC BY 4.0), distributed under the Creative Commons Attribution 4.0 International licence, as recorded in the arXiv licence metadata for this exact version and shown on the abstract page https://arxiv.org/abs/2508.20009v1. Full rights notice: licenses/arxiv-2508.20009-CC-BY-4.0.txt.\par\smallskip

\noindent\textit{Editorial scope.} Counted unit: the Proposition whose source label is \texttt{prop:QP-in-parallelotope} in the article (source0.tex lines 1468--1470, source label \texttt{prop:QP-in-parallelotope}), delivered together with the article's own complete original proof of it (source0.tex lines 1471--1520, one \texttt{proof} environment of 50 lines). No mathematics is written, repaired or re-derived: statement and proof are the article's own text line for line. Required context retained uncounted: lem:nvol (lines 485--491). Every definition, display and label of those lines is retained unchanged. Omitted from this package and not redistributed: 1--484, 492--1467, 1521--1968. Nothing inside a retained excerpt is omitted or altered: every retained body is delivered exactly as the article's lines are written, so no omission receipt of any kind accompanies this package. Citations: the retained excerpts carry no citation command at all, so no bibliography entry of the article is transcribed and no reference is redistributed. Reference adaptation: none was needed and none was made; every reference inside the delivered unit resolves inside it, so no unresolved reference or citation is produced. Printed slips kept literal: no printed word, symbol, hypothesis or number of the counted statement or of its proof was altered. Layout and class: the article's own class is \texttt{article} with packages including style; the wrapper is the lane's pinned \texttt{amsart} editorial wrapper at 10pt on A4, which restates the theorem-like environments the retained text needs and adds the article's own macro definitions that the retained text needs (\texttt{\textbackslash LD}, \texttt{\textbackslash Lcal}, \texttt{\textbackslash N}, \texttt{\textbackslash R}, \texttt{\textbackslash Z}, \texttt{\textbackslash ba}, \texttt{\textbackslash bb}, \texttt{\textbackslash blocks}, \texttt{\textbackslash bu}, \texttt{\textbackslash conv}, \texttt{\textbackslash lp}, \texttt{\textbackslash nvol}, \texttt{\textbackslash rem}, \texttt{\textbackslash rp}, \texttt{\textbackslash spn}, \texttt{\textbackslash st}), with extra wrapper packages (bm, bbm, dsfont, xspace, cancel, mathtools, cleveref). The delivered numbering is the unit's own. Assets: no retained span contains a figure environment, a graphics-inclusion command, a PDF-page inclusion, a table or a TikZ picture; no image file is copied into this package, the package declares no asset dependency and no image file is redistributed with it. Licence: version arXiv:2508.20009v1 of this article is distributed under the Creative Commons Attribution 4.0 International licence, as recorded in the arXiv licence metadata for this exact version and shown on the abstract page https://arxiv.org/abs/2508.20009v1. A full rights notice is delivered at licenses/arxiv-2508.20009-CC-BY-4.0.txt. Characters: every retained excerpt is pure ASCII, so the delivered engine, XeLaTeX with its default Latin Modern text font, reports no missing character.
\begin{lemma}\label{lem:nvol}
    Let $P$ be a lattice polytope and $L$ be a line that intersects $P$. Then
    \begin{align*}
    \left \lfloor \nvol(L \cap P) \right \rfloor \leq |L \cap P \cap \Z^d| \leq  \left \lfloor \nvol(L \cap P) \right \rfloor + 1
    \end{align*}
and the upper bound is attained if $L$ contains a lattice point on the boundary of $P$.
\end{lemma}

\begin{proposition}\label{prop:QP-in-parallelotope}
    Let $R$ be a rational parallelogram with parallel edges in direction $\bu$ that each contain an integral vertex. Then there exists a positive integer $q$ such that \mbox{$\LD_{R,u}:\N_{\geq q} \to \N$} agrees with a quasi-polynomial function of degree one whose period divides $q$.
\end{proposition}

\begin{proof}
    After applying a unimodular transformation, we may assume that $(0,0)$ is a vertex of $R$ and that $\bu = (1,0)$. Let the other two parallel edges, adjacent to those in direction $\bu$, have direction $(a,q)$ for some $a,q \in \Z, q >0$. Then, one of the edges in direction $\bu$ has the form $[(0,0),(\frac{p}{q},0)]$ and without loss of generality $p>0$. This edge contains a $\bu$-lattice diameter segment. Furthermore, the other edge emanating from the origin has the form $[(0,0),\lambda(a,q)]$ for some $\lambda = \frac{w}{q}$ and again we can assume $w \in \N$. That is, $\lambda(a,q) = (\frac{aw}{q},w)$.
    
    Let $k \in \N$ be a dilation factor. The lattice width of $kR$ in direction $(0,1)$ is $kw$.
    %The normalized length of lattice segments in direction $(0,1)$ agrees with their Euclidean length, since $\|(0,1)\|=1$, but we continue to write $\nvol$, because the formulas hold in general for the normalized length.
    For $j \in \{0, \ldots, kw\}$, define the lines on height $j$ as $L_j = (0,j) + \spn(\{(1,0)\})$ and denote the corresponding segments by $S_j(k) = L_j \cap kR$. When $k=1$ we just write $S_j$. Let $\Lcal(k) = \{S_j(k) \st j \in \{0, \ldots, kw\}\}$, that is, all $\bu$-lattice segments that intersect $kR$. For a segment $S \subseteq \R^2$, let $S_{{\Z}} := \conv (\{S \cap \Z^2\})$. So $S_j(k)$ contains a $\bu$-lattice diameter segment if and only if $S_j(k)_{\Z}$ is a $\bu$-lattice diameter segment.

    \emph{Claim 1:} $\nvol(S_j(k))-\nvol(S_{j}(k)_{\Z}) <1$ if and only if $S_{j}(k)_\Z$ is a $\bu$-lattice diameter segment.
    
    It suffices to show this for $k=1$ since $R$ is arbitrary. All $S_j$ have the same length, and $S_0$ is a lattice diameter segment since it contains a vertex. By \Cref{lem:nvol}
    \begin{align*}
        |S_0 \cap \Z^2| = \lfloor \nvol(S_0) \rfloor +1 
        = \lfloor \nvol(S_j) \rfloor +1  \geq |S_j \cap \Z^2| = \nvol({S_{j}}_{\Z}) +1.
    \end{align*}
    Thus $|S_0 \cap \Z^2| = |S_j \cap \Z^2|$ if and only if $\lfloor \nvol(S_j) \rfloor +1  = \nvol({S_{j}}_{\Z}) +1$ which holds if and only if $\nvol(S_j)- \nvol(S_{j{\Z}}) < 1$. This proves \textit{Claim 1}.


    \textit{Claim 2:} $\bu$-lattice diameter segments $S_j(k)_{\Z}$ are determined by $j \pmod q$. 
    
    Let $\chi_k: \{0, \ldots, kw\} \to \{0,1\}$ be the indicator function with $\chi_k(j)=1$ if and only if $L_j$ is a $\bu$-lattice diameter line of $kR$. Let $j,j' \in \{0, \ldots, kw\}$ with $j'= j+tq$ for some $t \in \N$. Then $S_{j'}(k) =S_j(k)+(ta,tq)$, and $S_{j'}(k)_\Z =S_j(k)_{\Z}+(ta,tq)$. Thus
    \begin{align*}
         \nvol(S_j(k))-\nvol(S_{j}(k)_{\Z})=\nvol(S_{j'}(k))-\nvol(S_{j'}(k)_{\Z})
    \end{align*}
    and by \textit{Claim 1} $S_j(k)_\Z$ is a $\bu$-lattice diameter segment of $kR$ if and only if $S_{j'}(k)_\Z$ is.
    
    \textit{Claim 3:} $\bu$-lattice diameter segments $S_j(k)_\Z$ are determined by $k \pmod q$.
    
    Let $j \in \{0, \ldots, kw\}$ and $i \in \{0, \ldots, q-1\}$. If $S_j(k) = [\ba,\bb]$, then $S_j(i) = [\ba,\bb + ((k-i)\frac{p}{q},0)]$. When $k \equiv i \pmod q$ the right endpoints of $S_j(k)$ and $S_j(i)$ differ by a vector in $\Z \times \{0\}$, so
    \begin{align*}
         \nvol(S_j(k))-\nvol(S_j(k)_\Z)=\nvol(S_j(i))-\nvol(S_j(i)_\Z).
    \end{align*}
    By \textit{Claim 1} it follows that $S_j(k)_\Z$ is a $\bu$-lattice diameter segment of $kR$ if and only if $S_j(i)_\Z$ is a $\bu$-lattice diameter segment of $iR$.
    
    \textit{The counting function:} Note that $q$ and $w$ are values that only depend on $R$. Let $i \in \{0, \ldots, q-1\}$, and let $k=\tilde{k}q+i$ for some $\tilde{k} \in \N$. Define a $q$-block of $kR$ as the set of lines $L_{tq}, \ldots, L_{tq+q-1}$ for those $t \in \N$ for which $tq+q-1 \leq kw$. The number of $q$-blocks in $kR$ is
    \begin{align*}
         \blocks_i(k)  = \left \lfloor  \frac{kw+1}{q} \right\rfloor = \tilde{k}w + \left\lfloor \frac{iw+1}{q}\right\rfloor = \frac{wk}{q}  - \lp \frac{iw}{q} - \left\lfloor \frac{iw+1}{q} \right\rfloor \rp,
    \end{align*}
    which is a linear function in $k$. Next we count the number of $\bu$-lattice diameter lines in $kR$ per $q$-block. By \textit{Claim 3}, this only depends on the congruence class $k \bmod q$. Take $q+i$ as a representative for $i \bmod q$,
    % The width of $(q+i)R$ in direction $(0,1)$ is $(q+i)w>q$, so $(q+i)R$ 
    since it has at least one $q$-block, and let
    \begin{align*}
        n_i := \sum_{j=0}^{q-1}\chi_{q+i}(j) = \# \{  L_j: j =0, \ldots, q-1,\, L_j \text{ is a $\bu$-lattice diameter line of } (q+i)P\}.
    \end{align*}
    Define $\rem_i$ as the number of lattice lines intersecting $kR$ not contained in a $q$-block, then $\rem_i= iw+1 \bmod q$. Clearly $kw +1 \equiv iw+1 \pmod q$. 
    The number of those lines that are $\bu$-lattice diameter lines is
    \begin{align*}
        r_i := \sum_{j=0}^{\rem_i} \chi_{i}(j)  = \# \{  L_j: j =0, \ldots, \rem_i, \, L_j \text{ is a $\bu$-lattice diameter line of } iP\}.
    \end{align*}
    Together we get $\LD_{R,u\,}(k) = n_i \cdot \blocks_{i}(k) + r_i$ which for fixed congruence class $ i \bmod q$ is a linear function in $k$. Thus $\LD_{R,u}$ agrees with a quasi-polynomial function of degree one, whose period \mbox{divides $q$}.
\end{proof}

\end{document}
